\section{Related Work}

CNT sits near several existing lines of work, but its central claim is different. The nearest contrast is with reasoning methods that improve performance without explicitly isolating a local necessity object. Chain-of-thought prompting and self-consistency improve reasoning behavior at inference time, but they do not by themselves identify which local steps are causally or conditionally necessary for later success \citep{wei2023chainofthoughtpromptingelicitsreasoning,wang2023selfconsistencyimproveschainthought}.

The second nearby line is process supervision. Process- and outcome-based feedback for math reasoning, as well as large verifier-style datasets, aim to score or rank intermediate reasoning steps \citep{uesato2022solvingmathwordproblems,lightman2023letsverifystepstep}. CNT is compatible with this literature, but it targets a narrower and more audited object: not generic step quality, but trajectory-conditional necessity under controlled edits and continuators. This is why the current project emphasizes synthetic recovery, cross-editor invariance, held-out continuator stability, and deletion-paraphrase asymmetry before making strong training claims.

The third nearby line is preference optimization. Standard pairwise preference objectives such as DPO provide a flexible way to train from comparisons \citep{rafailov2024directpreferenceoptimizationlanguage}. Our Stage~B experiments borrow that basic optimization shape, but the current results suggest that preference-style offline objectives are not automatically enough. Even when several structurally different families improve offline margins, the rollout utility signal can remain tied or fall onto a utility-versus-weighted-necessity frontier. This paper therefore uses preference learning as an instrument for testing the object, not as a guarantee that the object will immediately convert into utility gains.
